\documentclass[a4paper,12pt]{article}
\usepackage{amsmath,amssymb,graphicx,cite}
\usepackage[utf8]{inputenc}
\usepackage[colorlinks=true,linkcolor=blue,citecolor=blue]{hyperref}

\title{Supplementary Material: Subharmonic Corrections in Gravitational Wave Ringdowns}
\author{Haobo Ma\textsuperscript{1} \and Wen Niu\textsuperscript{1}}
\date{\vspace{-5ex}}

\begin{document}

\maketitle

\section{Mathematical Derivations}
Here we provide detailed mathematical derivations for the subharmonic extension to the quasi-normal mode model.

\subsection{Extended QNM Model}
The standard QNM model is extended with a subharmonic correction term:
\begin{equation}
h(t) = \sum_n A_n e^{-i\omega_n t} e^{-t/\tau_n}[1 + \varepsilon \sin((1+\varepsilon)\omega_n t)]
\end{equation}

\section{Resonance Optimization Algorithm}
To detect the optimal value of the subharmonic parameter $\varepsilon$, we employ a resonance optimization algorithm:

\begin{enumerate}
\item Initialize parameter search space for $\varepsilon \in [0.1, 0.3]$
\item For each test value $\varepsilon_i$:
   \begin{enumerate}
   \item Compute extended waveform $h(t;\varepsilon_i)$
   \item Calculate goodness-of-fit measure $\chi^2(\varepsilon_i)$
   \item Determine resonance strength $R(\varepsilon_i)$
   \end{enumerate}
\item Identify optimal $\varepsilon^*$ maximizing resonance strength and minimizing $\chi^2$
\end{enumerate}

\section{Connection to Fundamental Constants}
The observed value $\varepsilon \approx 0.18$ has intriguing connections to fundamental constants:

\begin{align}
\varepsilon &\approx \frac{\phi-1}{10} \approx 0.1618\\
\varepsilon &\approx 25\alpha \approx 0.1825\\
\varepsilon &\approx \sqrt{\frac{\ell_P}{R_S}} \cdot \frac{1}{\pi} \approx 0.179
\end{align}

where $\phi$ is the golden ratio, $\alpha$ is the fine structure constant, $\ell_P$ is the Planck length, and $R_S$ is the Schwarzschild radius.

\section{Additional Statistical Tests}
We performed extensive statistical tests to validate the significance of the subharmonic parameter across multiple black hole merger events, confirming its robustness against variations in black hole mass, spin, and merger dynamics.

\bibliographystyle{unsrt}
\bibliography{references}

\end{document} 